\section{Spécifications techniques}
\label{sec:specs-tech}

\subsection{Choix technologiques}

\begin{table}[H]
    \centering
    \begin{tabularx}{\textwidth}{lX}
        \toprule
        \textbf{Couche} & \textbf{Choix retenu et justification}\\
        \midrule
        Langage backend & \textbf{Python~3.11+} --- écosystème data/ML
        mature, lisibilité, alignement avec les compétences enseignées.\\
        Framework web & \textbf{Django~5.x} --- ORM puissant, admin
        gratuit, support i18n natif via \texttt{gettext}, mature.\\
        Base de données & \textbf{SQLite} en développement,
        \textbf{PostgreSQL} en production --- migration transparente via
        l'ORM Django.\\
        Frontend & \textbf{Templates Django + Tailwind CSS~3} --- rendu
        serveur compatible \texttt{\{\% trans \%\}}, design moderne sans
        framework JS lourd.\\
        Interactivité & \textbf{HTMX + Alpine.js} --- drag-and-drop Kanban,
        widget chatbot, sans bundler ni étape de build JS.\\
        IA conversationnelle & \textbf{API DeepSeek}
        (\texttt{deepseek-chat}) --- API compatible OpenAI, faible coût,
        bonne qualité multilingue.\\
        Apprentissage automatique & \textbf{scikit-learn}
        (\texttt{RandomForestClassifier}) --- robuste, interprétable,
        adapté à un petit jeu de données synthétique.\\
        Sérialisation modèle & \textbf{joblib} --- standard de fait pour
        sklearn.\\
        Configuration & \textbf{python-decouple} + fichier \texttt{.env}
        --- séparation code/secrets.\\
        Tests & \textbf{pytest-django} + \texttt{coverage.py}.\\
        \bottomrule
    \end{tabularx}
    \caption{Stack technique retenue}
    \label{tab:stack}
\end{table}

\subsection{Exigences non fonctionnelles}

\begin{description}
    \item[Performance] Toutes les pages doivent répondre en moins de
    300~ms sur un dataset de démonstration. Les requêtes du tableau de
    bord ROI utilisent \texttt{annotate()} et \texttt{aggregate()} pour
    rester en une seule requête SQL.
    \item[Sécurité] Protection CSRF native Django, requêtes sortantes
    DeepSeek effectuées exclusivement côté serveur (la clé API n'est
    jamais exposée au navigateur), validation des uploads de créatifs
    (taille max, MIME type whitelist).
    \item[Maintenabilité] Découpage en applications Django distinctes
    (\texttt{pipeline}, \texttt{assets}, \texttt{finance}, \texttt{chatbot},
    \texttt{prediction}, \texttt{core}), couverture de tests $\geq 60~\%$
    sur la logique métier.
    \item[Portabilité] Aucune dépendance native non-pip~; le projet doit
    démarrer sous macOS, Linux et Windows en moins de 5~minutes.
    \item[Internationalisation] Toutes les chaînes user-facing passent par
    \texttt{gettext\_lazy} ou \texttt{\{\% trans \%\}}. Les fichiers
    \texttt{.po} sont versionnés~; la compilation \texttt{.mo} est
    documentée dans le \texttt{README}.
\end{description}

\subsection{Variables d'environnement}

\begin{lstlisting}[language=bash,caption={Fichier .env.example},label=lst:env]
DJANGO_SECRET_KEY=change-me-in-prod
DJANGO_DEBUG=True
DEEPSEEK_API_KEY=sk-...
DEEPSEEK_BASE_URL=https://api.deepseek.com/v1
DEEPSEEK_MODEL=deepseek-chat
\end{lstlisting}
